% ============================================================================
% SECTION 9: CONCLUSION
% ============================================================================
\section{Conclusion}\label{sec:conclusion}

We presented a boundary information-geometric framework for the local S-box
distributions that arise in side-channel key recovery.  The main modelling
point is the separation between three spaces: the regular full-support Walsh
exponential family, the singular key-induced boundary orbit, and the
adversary's belief simplex.  This separation makes Fisher, KL, and projection
statements mathematically well scoped.

For bijective S-boxes, Walsh coordinates give closed-form boundary covariance
matrices.  The active covariance subspace has dimension $2^s-1$ and all active
eigenvalues equal $2^s$ in the unnormalised Walsh basis.  Thus this part of the
geometry is a general structural consequence of bijectivity.

For CPA, Bayesian updates occur on the belief simplex.  Under Hamming-weight
Gaussian leakage, the scalar variance information is a baseline fixed by
bijectivity, while pairwise leakage gaps determine expected posterior-odds
growth against wrong keys.  Computations on PRESENT, GIFT, and SKINNY
illustrate this distinction and reproduce the active curvature diagnostic
$K=-1/4$ in the tested 4-bit cases.

The framework is not a new S-box classification and not a standalone attack.
Its contribution is a controlled geometric language for connecting singular
S-box distributions, boundary covariance, and side-channel Bayesian inference.
The security value of the framework is to make explicit which geometric
quantities are structural invariants and which quantities affect local
key-hypothesis discrimination under a specified leakage model.
